\clearpage
\begin{proofbox}
\textbf{\textcolor{CProof}{证明一｜最快：点差法三步速写（考场首选）}}
\begin{center}\begin{tikzpicture}[x=1.45cm,y=1.45cm,>=stealth,font=\small]
\node[font=\small\bfseries] at (0,3.15) {点差法：坐标差给方向，坐标和给中点};
\draw[->,gray] (-3.5,0)--(3.7,0) node[right] {$x$};
\draw[->,gray] (0,-2.4)--(0,2.5) node[above] {$y$};
\draw[CStatement,thick] (0,0) ellipse (3 and 2);
\coordinate (A) at (3,0); \coordinate (B) at (0,2); \coordinate (M) at (1.5,1);
\draw[teal,very thick] (A)--(B);
\fill (A) circle(1.8pt) node[below right] {$A$};
\fill (B) circle(1.8pt) node[above left] {$B$};
\fill[orange!80!black] (M) circle(2pt) node[above right] {$M$};
\node[below left] at (0,0) {$O$};
\draw[purple,dashed] (B)--(3,2)--(A);
\node[above] at (1.5,2) {$x_1-x_2$};
\node[right] at (3,1.25) {$y_2-y_1$};
\node[below,align=center] at (0,-2.85) {$A(x_1,y_1),\ B(x_2,y_2),\ M\bigl(\frac{x_1+x_2}{2},\frac{y_1+y_2}{2}\bigr)$};

\end{tikzpicture}
\end{center}
设 $A(x_1,y_1),B(x_2,y_2)$ 是不同端点，$M=(x_0,y_0)$。
\begin{enumerate}[leftmargin=*,itemsep=0.25em]
\item 两个椭圆方程相减，并用平方差：
\[ \frac{(x_1-x_2)(x_1+x_2)}{a^2}+
      \frac{(y_1-y_2)(y_1+y_2)}{b^2}=0.\]
\item 用 $x_1+x_2=2x_0,\ y_1+y_2=2y_0$：
\[ \boxed{\frac{x_0}{a^2}(x_1-x_2)+\frac{y_0}{b^2}(y_1-y_2)=0.}\]
\item 若 $y_0\ne0$，则必有 $x_1\ne x_2$（否则上式迫使 $y_1=y_2$，与端点不同矛盾）。故
\[\frac{x_0}{a^2}+\frac{y_0}{b^2}\underbrace{\frac{y_1-y_2}{x_1-x_2}}_{k}=0
\ \Longrightarrow\ \boxed{k=-\frac{b^2x_0}{a^2y_0}}.\]
\end{enumerate}
\textbf{为何最快：}只用“平方差 $\to$ 中点和 $\to$ 斜率”，不用设直线与椭圆联立，更不用求交点根。
\end{proofbox}
\clearpage
\begin{proofbox}
\textbf{\textcolor{CProof}{证明二｜最低门槛：中点正负位移法（只引入一个位移）}}
\begin{center}\begin{tikzpicture}[x=1.45cm,y=1.45cm,>=stealth,font=\small]
\node[font=\small\bfseries] at (0,3.15) {位移法：从中点沿弦作相反位移};
\draw[->,gray] (-3.5,0)--(3.7,0) node[right] {$x$};
\draw[->,gray] (0,-2.4)--(0,2.5) node[above] {$y$};
\draw[CStatement,thick] (0,0) ellipse (3 and 2);
\coordinate (A) at (3,0); \coordinate (B) at (0,2); \coordinate (M) at (1.5,1);
\draw[teal,very thick] (A)--(B);
\fill (A) circle(1.8pt) node[below right] {$A$};
\fill (B) circle(1.8pt) node[above left] {$B$};
\fill[orange!80!black] (M) circle(2pt) node[above right] {$M$};
\node[below left] at (0,0) {$O$};
\draw[orange!80!black,->,thick] (1.5,.83)--(2.9,-.10);
\draw[purple,->,thick] (1.5,1.17)--(.1,2.10);
\node[below,fill=white,inner sep=2pt] at (2,-.65) {$(t,kt)$};
\node[above,fill=white,inner sep=2pt] at (-.8,1.25) {$(-t,-kt)$};
\node[below,align=center] at (0,-2.85) {$A=M+(t,kt),\qquad B=M-(t,kt)$};

\end{tikzpicture}
\end{center}
先考虑斜率存在的弦。设 $A(x_0+t,y_0+kt)$、$B(x_0-t,y_0-kt)$，其中 $t\ne0$。\textbf{为什么这样设：}“中点”就是两端点关于 $M$ 对称，横坐标是 $x_0\pm t$；同一条直线的纵坐标变化为 $k(\pm t)$。
把 $A,B$ 代入椭圆并相减：
\[
\frac{(x_0+t)^2-(x_0-t)^2}{a^2}+
\frac{(y_0+kt)^2-(y_0-kt)^2}{b^2}=0.
\]
利用 $(u+v)^2-(u-v)^2=4uv$，直接得到
\[
4t\left(\frac{x_0}{a^2}+\frac{k y_0}{b^2}\right)=0
\ \Longrightarrow\ \boxed{k=-\frac{b^2x_0}{a^2y_0}}\quad(y_0\ne0).
\]
\textbf{优点：}仅需理解“沿弦向两侧等距离”，无需记 $x_1+x_2$ 之类中间变量；\textbf{局限：}此设法默认斜率存在，竖直弦应使用证明一的无除法式。
\end{proofbox}
\clearpage
\begin{proofbox}
\textbf{\textcolor{CProof}{证明三｜常规法：端点方程、平方差、中点坐标（完整书写）}}
\begin{center}\begin{tikzpicture}[x=.58cm,y=.58cm,font=\small,>=stealth]
\foreach \s in {0,7.6,15.2}{
\begin{scope}[xshift=\s*0.58cm]
\draw[gray,->] (-3.3,0)--(3.5,0);
\draw[gray,->] (0,-2.3)--(0,2.5);
\draw[CStatement,thick] (0,0) ellipse(3 and 2);
\end{scope}}
\draw[teal,very thick] (-2.598076,1)--(2.598076,1);
\fill[orange!80!black] (0,1) circle(2pt) node[above right] {$M$};
\node at (0,3.2) {$x_0=0,\ y_0\ne0$}; \node at (0,-3) {水平：$y=y_0$};
\begin{scope}[xshift=4.408cm]
\draw[teal,very thick] (1.5,-1.732051)--(1.5,1.732051);
\fill[orange!80!black] (1.5,0) circle(2pt) node[above right] {$M$};
\node at (0,3.2) {$y_0=0,\ x_0\ne0$}; \node at (0,-3) {竖直：$x=x_0$};
\end{scope}
\begin{scope}[xshift=8.816cm]
\draw[teal,very thick] (-3,0)--(3,0);
\draw[purple,dashed,thick] (-2.12132,-1.414214)--(2.12132,1.414214);
\fill[orange!80!black] (0,0) circle(2pt) node[above left] {$M=O$};
\node at (0,3.2) {$x_0=y_0=0$}; \node at (0,-3) {中心弦：方向任意};
\end{scope}
\end{tikzpicture}
\end{center}
\begin{enumerate}[leftmargin=*,itemsep=0.35em]
\item \textbf{设点并写中点。} $A(x_1,y_1),B(x_2,y_2)$，
$2x_0=x_1+x_2,\ 2y_0=y_1+y_2$。
\item \textbf{分别代入。}
$\frac{x_1^2}{a^2}+\frac{y_1^2}{b^2}=1$，
$\frac{x_2^2}{a^2}+\frac{y_2^2}{b^2}=1$。
\item \textbf{两式相减分解。}
$\frac{(x_1-x_2)(x_1+x_2)}{a^2}+
\frac{(y_1-y_2)(y_1+y_2)}{b^2}=0$。
\item \textbf{代回中点。} $\frac{x_0}{a^2}(x_1-x_2)+\frac{y_0}{b^2}(y_1-y_2)=0$。
\item \textbf{分情形整理。} 若 $y_0\ne0$，除以 $x_1-x_2\ne0$，得 $k=-b^2x_0/(a^2y_0)$。若 $y_0=0,x_0\ne0$，必有 $x_1=x_2$，是竖直弦；若 $x_0=0,y_0\ne0$，必有 $y_1=y_2$，是水平弦；若 $M=O$，方向不唯一。
\end{enumerate}
\textbf{为何保留常规法：}完整交代除法依据和所有边界，适用于书面证明、讲授和自查。
\end{proofbox}
\clearpage
\begin{proofbox}
\textbf{\textcolor{CProof}{证明四｜补充：椭圆拉伸成圆（直观几何法）}}
\begin{center}\begin{tikzpicture}[x=1cm,y=1cm,font=\small,>=stealth]
\node at (0,3) {原椭圆：$a=3,\ b=2$};
\draw[CStatement,thick] (0,0) ellipse (3 and 2);
\draw[teal,very thick] (3,0)--(0,2);
\draw[orange!80!black,dashed,thick] (0,0)--(1.5,1);
\fill (3,0) circle(1.7pt) node[right] {$A$};
\fill (0,2) circle(1.7pt) node[above] {$B$};
\fill (1.5,1) circle(1.7pt) node[above right] {$M$};
\fill (0,0) circle(1.7pt) node[below] {$O$};
\draw[->,thick] (3.7,.4)--(5.2,.4) node[midway,above] {$X=x/a$} node[midway,below] {$Y=y/b$};
\begin{scope}[xshift=8cm]
\node at (0,3) {单位圆（放大显示）};
\draw[CStatement,thick] (0,0) circle(2);
\draw[teal,very thick] (2,0)--(0,2);
\draw[orange!80!black,dashed,thick] (0,0)--(1,1);
\draw (1.15,.85)--(1,.7)--(.85,.85);
\fill (2,0) circle(1.7pt) node[right] {$A'$};
\fill (0,2) circle(1.7pt) node[above] {$B'$};
\fill (1,1) circle(1.7pt); \node at (.1,1.2) {$M'$};
\fill (0,0) circle(1.7pt) node[below] {$O$};
\node at (0,-2.7) {$OM'\perp A'B'$};
\end{scope}
\node at (0,-2.7) {伸缩一般不保角};
\end{tikzpicture}
\end{center}
坐标变换 $X=x/a,\ Y=y/b$ 将椭圆化为单位圆。设变换后弦端点为 $A',B'$，中点为 $M'=(x_0/a,y_0/b)$。圆心到非中心弦中点的向量与弦方向垂直（中心弦时该向量为零，内积仍为零），故
\[
\bigg(\frac{x_0}{a},\frac{y_0}{b}\bigg)\cdot
\bigg(\frac{x_1-x_2}{a},\frac{y_1-y_2}{b}\bigg)=0.
\]
即 $\frac{x_0}{a^2}\Delta x+\frac{y_0}{b^2}\Delta y=0$，以下同证明一。
\textbf{为什么值得补充：}它解释了“为什么是负号”和“为什么是 $b^2/a^2$”，便于理解、迁移和长期记忆；但在考场上通常慢于点差速写。
\end{proofbox}
\clearpage
\begin{proofbox}
\textbf{\textcolor{CProof}{由必要到充分：为何内部非中心点一定对应一条真实弦？}}
\begin{center}\begin{tikzpicture}[x=1.45cm,y=1.45cm,>=stealth,font=\small]
\node[font=\small\bfseries] at (0,3.15) {真实弦：中点在内部；弦线只与小椭圆相切};
\draw[->,gray] (-3.5,0)--(3.7,0) node[right] {$x$};
\draw[->,gray] (0,-2.4)--(0,2.5) node[above] {$y$};
\draw[CStatement,thick] (0,0) ellipse (3 and 2);
\coordinate (A) at (3,0); \coordinate (B) at (0,2); \coordinate (M) at (1.5,1);
\draw[teal,very thick] (A)--(B);
\fill (A) circle(1.8pt) node[below right] {$A$};
\fill (B) circle(1.8pt) node[above left] {$B$};
\fill[orange!80!black] (M) circle(2pt) node[above right] {$M$};
\node[below left] at (0,0) {$O$};
\draw[purple,dashed,thick] (0,0) ellipse (2.121320344 and 1.414213562);
\draw[orange!80!black,dashed] (0,0)--(M);
\node[align=center] at (-.8,-.8) {同心小椭圆\\$r^2=1/2$};
\node[below,align=center] at (0,-2.85) {$M\pm s(1,-2/3),\ s=3/2$：两交点关于 $M$ 对称};
\end{tikzpicture}
\end{center}
设 $0<r^2=\frac{x_0^2}{a^2}+\frac{y_0^2}{b^2}<1$。过 $M$ 取方向向量 $v=(u,v_y)\ne(0,0)$，令它满足
$\frac{x_0}{a^2}u+\frac{y_0}{b^2}v_y=0$（这样的方向唯一，差非零倍数）。
直线上的点写成 $(x,y)=(x_0,y_0)+t(u,v_y)$。代入椭圆：
\[
 r^2+t^2\bigg(\frac{u^2}{a^2}+\frac{v_y^2}{b^2}\bigg)=1
 \quad\Longrightarrow\quad
 t=\pm\sqrt{\frac{1-r^2}{u^2/a^2+v_y^2/b^2}}.
\]
两根互为相反数且均非零，所以确实是以 $M$ 为中点的一条\textbf{非退化弦}；若 $M=O$，每个方向都满足垂直条件，故不唯一。\textbf{这也说明：}中点在内部是弦存在性的几何条件，并不是代数“相减”步骤额外要求的条件。
\end{proofbox}
